\newpage
\section{FEASIBILITY ANALYSIS}

\subsection{Technical Feasibility}
The proposed Audio Deepfake Detection (ADD) system is technically feasible due to the availability of public datasets, deep learning frameworks, and audio processing tools. The system can be developed using Python, PyTorch, Librosa, and related libraries. Existing RTC technologies and audio simulation tools can be used to emulate real-world VoIP conditions such as codec compression, packet loss, and network delay. The required computational resources are available through local GPU systems or cloud platforms.

\subsection{Operational Feasibility}
The system is operationally feasible because it addresses a practical security concern in modern RTC and VoIP platforms. The proposed detector operates as an external service and can be integrated into existing communication systems without modifying their core functionality. This allows real-time monitoring while preserving normal communication workflows.

\subsection{Economic Feasibility}
The project requires minimal financial investment since most development tools, frameworks, and datasets are open source. The primary costs involve computing resources, internet access, and optional cloud infrastructure for model training and deployment. These costs are manageable within the scope of a student project.

\subsection{Schedule Feasibility}
The project can be completed within the allocated academic timeframe. Development activities including literature review, dataset preparation, model development, RTC simulation, system integration, evaluation, and documentation can be planned and executed incrementally throughout the project duration.

\subsection{Ethical and Legal Feasibility}
The project uses publicly available datasets and focuses on detecting synthetic speech for security and research purposes. No private voice data will be collected without consent. The proposed system is intended as a defensive tool to improve trust and security in communication systems and does not interfere with normal user communications.